\section{Discussion and Limitations}\label{sec:discussion}

Table~\ref{tab:noise_sensitivity} evaluates fixed-threshold sensitivity under Gaussian perturbation. The ranking metrics degrade gradually under mild perturbation, while fixed-threshold decision quality becomes more sensitive under stronger perturbation; this supports a conservative deployment recommendation with drift monitoring and periodic recalibration.

\begin{table*}[t]
  \centering
  \caption{Fixed-threshold sensitivity under Gaussian input perturbation. The threshold \(\tau=0.425420\) is calibrated on the clean independent benign split and then kept fixed for all perturbed test evaluations.}
  \label{tab:noise_sensitivity}
  \small
  \resizebox{\textwidth}{!}{%
  \begin{tabular}{ccccccc}
    \toprule
    Perturbation \(\sigma\) & AUC & AP & Test FPR@\(\tau\) & Precision@\(\tau\) & Recall@\(\tau\) & F1@\(\tau\) \\
    \midrule
    0.0000 & 0.9783 & 0.9631 & 0.0460 & 0.9384 & 0.9568 & 0.9475 \\
    0.0100 & 0.9778 & 0.9613 & 0.0724 & 0.9078 & 0.9739 & 0.9397 \\
    0.0300 & 0.9733 & 0.9547 & 0.3520 & 0.6738 & 0.9925 & 0.8027 \\
    0.0500 & 0.9647 & 0.9428 & 0.9636 & 0.4319 & 0.9998 & 0.6032 \\
    \bottomrule
  \end{tabular}
  }
\end{table*}


The calibration implication is practical. A threshold that is well calibrated on one benign regime can drift under changed traffic conditions even when score ranking remains moderately stable. For deployment, this suggests a two-stage policy: monitor benign alarm rate continuously, and trigger recalibration when observed benign FPR exceeds an operational tolerance band for a sustained period.

Current evidence is intentionally compact and protocol-consistent. CICIDS2017 is the primary evidence chain, while SWaT and UNSW-NB15 are used to characterize operating boundaries under additional conditions rather than as co-equal primary benchmarks.

The study has four key limitations. First, Bot recall is zero at the selected CICIDS2017 operating point, showing attack-family coverage gaps despite strong aggregate metrics. Second, fixed-threshold behavior is sensitive under stronger perturbation, and UNSW-NB15 further exposes threshold-transfer failure under distribution shift. Third, the baseline suite is representative under strict protocol alignment, not exhaustive across all modern architectures. Fourth, audit evidence is lightweight and empirical: attribution-guided masking supports alarm-trace inspection, but does not prove complete faithfulness.

\subsection{Ethical Considerations and Open-Science Notes}
This work uses public benchmark traffic data and does not involve active attacks on operational systems. The main risk is overextending benchmark conclusions to production networks. We mitigate this risk by reporting observed benign FPR, explicit failure modes (including Bot miss and UNSW transfer failure), and strict calibration provenance.

The artifact structure is designed for reproducibility: strict split settings, calibration protocol, selected checkpoints, and table-generation scripts are recorded in the project outputs. We encourage deployment as analyst-assistive triage, not as a standalone replacement for security investigation workflows.
